%&LaTeX
% $Id: ICandBC.tex,v 1.7 2003-05-21 23:49:03 cjkimme Exp $

\section{Initial conditions and boundary conditions}
\label{set.ICandBC}
Initial conditions and boundary conditions may be specified on 
selected nodes. Additionally, the variational form 
of governing equations leads to natural boundary conditions that may 
be specified on selected parts of the domain boundary. \tahoe 
supports a number of geometry database types. These are listed in 
Table~\ref{tab.input.file.type}. The database format determines how 
nodes or sets of nodes are identified and how parts of the domain 
boundary are identified.
\begin{table}[h]
\begin{center}
\caption{\label{tab.BC.type} Identification of nodes and 
boundaries.}
\begin{tabular}[c]{|c|c|c|}
\hline
 \parbox[b]{1.25in}{\centering \textbf{database}}
&\parbox[b]{1.5in}{\centering \textbf{node ID}}
&\parbox[b]{1.5in}{\centering \textbf{side ID}}\\
\hline
  \parbox[b]{1.25in}{\centering TahoeI}
& \parbox[b]{1.5in}{\centering node}
& \parbox[b]{1.5in}{\centering element,facet} \\
\hline
  \parbox[b]{1.25in}{\centering TahoeII}
& \parbox[b]{1.5in}{\centering node set ID}
& \parbox[b]{1.5in}{\centering side set ID} \\
\hline
  \parbox[b]{1.25in}{\centering \textsf{ExodusII}~\cite{ExodusII}}
& \parbox[b]{1.5in}{\centering node set ID}
& \parbox[b]{1.5in}{\centering side set ID} \\
\hline
\end{tabular}
\end{center}
\end{table}
Table~\ref{tab.BC.type} lists how nodes and element sides are 
identified for the different database types. Their use will be 
explained with examples below. For the TahoeI format, boundary 
segments are explicitly specified with an element and facet number. 
Canonical facet numbering for the different integration domain 
geometries supported by \tahoe are listed in 
Section~\ref{sect.integration.domains}. For TahoeII and 
\textsf{ExodusII} databases, groups of nodes may be collected into 
node sets and groups of element facets are collected into side sets. 
Both types of sets are identified with ID numbers.

Many of the boundary conditions may be specified as functions of 
time. Their values have the form
\begin{equation}
\label{eq.func.time}
    g\farg{t} = g\,f^{\rbrkt{i}}\farg{t},
\end{equation}
where $g$ is a scalar coefficient and $f^{\rbrkt{i}}\farg{t}$ is a
schedule function. Schedule functions are explained in 
Section~\ref{sect.sched.param}. Time varying functions in the input 
file are specified by providing $g$ and $i$, as defined 
in~\eqref{eq.func.time}.

\subsection{Nodal initial conditions}
\label{set.ICandBC.IC}
Each initial condition is specified as
\[
\sbrkt{\textit{node ID}} \quad 
\sbrkt{\textit{degree of freedom}} \quad
\sbrkt{\textit{order}} \quad
\sbrkt{\textit{value}}
\]
The definition of the \textit{node ID} depends on the geometry 
database type as listed in Table~\ref{tab.BC.type}. The \textit{order} refers to the time derivative of the field variable to which the initial conditions applies. A \textit{node ID} of \texttt{-1} applies the boundary condition to \textit{all} nodes.

\subsection{Nodal kinematic boundary conditions}
\label{set.ICandBC.KBC}
Each nodal kinematic boundary condition is specified as
\[
\sbrkt{\textit{node ID}} \quad 
\sbrkt{\textit{degree of freedom}} \quad
\sbrkt{\textit{code}} \quad
\sbrkt{\textit{schedule}} \quad
\sbrkt{\textit{value}} 
\]
The definition of the \textit{node ID} depends on the geometry 
database type as listed in Table~\ref{tab.BC.type}. The \textit{code} 
specifies if the condition is applied to the nodal unknown $u$, or one 
of its time derivatives $\dot{u}$, $\ddot{u}$, or higher. The codes 
are listed in Table~\ref{tab.KBC.code}.
\begin{table}[h]
\begin{center}
\caption{\label{tab.KBC.code} Codes for nodal kinematic boundary 
conditions.}
\begin{tabular}[c]{|c|c|}
\hline
 \parbox[b]{0.75in}{\centering \textbf{code}}
&\parbox[b]{0.75in}{\centering \textbf{type}}\\
\hline
 \parbox[b]{0.75in}{\centering 0}
&\parbox[b]{0.75in}{\centering \textit{fixed}}\\
\hline
 \parbox[b]{0.75in}{\centering 1}
&\parbox[b]{0.75in}{\centering $u$}\\
\hline
 \parbox[b]{0.75in}{\centering 2}
&\parbox[b]{0.75in}{\centering $\dot{u}$}\\
\hline
 \parbox[b]{0.75in}{\centering 3}
&\parbox[b]{0.75in}{\centering $\ddot{u}$}\\
\hline
\end{tabular}
\end{center}
\end{table}
The \textit{value} and \textit{schedule} in the nodal kinematic 
boundary condition specification refer to 
$g$ and $i$ in~\eqref{eq.func.time}, respectively, for the variation 
of the boundary condition in time.
For fields approximated by meshfree shape functions, which lack the Kronecker delta property,
a kinematic boundary condition prescribes the nodal coefficient rather than the physical value.
Adding \texttt{collocation="1"} to a \texttt{kinematic\_BC} entry prescribes the physical value
instead, see Section~\ref{set.ICandBC.meshfree.EBC}.

\subsection{Kinematic boundary condition controllers}
\label{set.ICandBC.KBC.controller}

\begin{table}[h]
\caption{\label{tab.KBC.controller.type} Kinematic boundary condition controllers.}
\begin{center}
{\small
\begin{tabular}[c]{|c|p{3.05in}|>{\raggedright\arraybackslash}p{2.45in}|}
\hline
\textbf{code} & \textbf{XML name} & \textbf{description}\\
\hline
0 & \texttt{K-field} & K-field\\
\hline
1 & \texttt{bi-material\_K-field} & bimaterial K-field\\
\hline
2 & \texttt{mapped\_nodes} & Mapped displacements\\
\hline
3 & \texttt{tied\_nodes} & Tied displacements\\
\hline
5 & \texttt{periodic\_nodes} & Periodic displacements\\
\hline
7 & \texttt{scaled\_velocity} & Scaled velocity\\
\hline
9 & \texttt{torsion} & Torsion\\
\hline
10 & \texttt{conveyor} & Conveyor (moving periodic domain)\\
\hline
11 & \texttt{symmetric\_conveyor} & Symmetric conveyor\\
\hline
12 & \texttt{collocation\_KBC} & Direct nodal collocation of essential boundary conditions for meshfree fields, Section~\ref{set.ICandBC.meshfree.EBC}\\
\hline
\end{tabular}}
\end{center}
\end{table}

\subsubsection{K-field}
\subsubsection{Interface K-field}
\subsubsection{Mapped displacements}
\subsubsection{Tied displacements}
\subsubsection{Periodic displacements}
\subsubsection{Scaled velocity}
\subsubsection{Torsion}
\subsubsection{Conveyor}
\subsubsection{Direct nodal collocation}
See Section~\ref{set.ICandBC.meshfree.EBC}.

\subsection{Applied nodal forces}
\label{set.ICandBC.FBC}
Each nodal force boundary condition is specified as
\[
\sbrkt{\textit{node ID}} \quad 
\sbrkt{\textit{degree of freedom}} \quad
\sbrkt{\textit{schedule}} \quad
\sbrkt{\textit{value}}
\]
The definition of the \textit{node ID} depends on the geometry 
database type as listed in Table~\ref{tab.BC.type}. 
The \textit{value} and \textit{schedule} in the nodal kinematic 
boundary condition specification refer to 
$g$ and $i$ in~\eqref{eq.func.time}, respectively, for the variation 
of the boundary condition in time.

\subsection{Force controllers}
\label{set.ICandBC.FBC.controller}

\begin{table}[h]
\caption{\label{tab.FBC.controller.type} Force boundary condition controllers.}
\begin{center}
{\small
\begin{tabular}[c]{|c|p{3.05in}|>{\raggedright\arraybackslash}p{2.45in}|}
\hline
\textbf{code} & \textbf{XML name} & \textbf{description}\\
\hline
0 & \texttt{wall\_penalty} & Rigid barrier (penalty)\\
\hline
1 & \texttt{sphere\_penalty} & Rigid sphere (penalty)\\
\hline
2 & \texttt{sphere\_augmented\_Lagrangian} & Rigid sphere (augmented Lagrangian)\\
\hline
3 & \texttt{sphere\_penalty\_meshfree} & Rigid sphere for meshfree domains (penalty)\\
\hline
4 & \texttt{wall\_augmented\_Lagrangian} & Rigid barrier (augmented Lagrangian)\\
\hline
5 & \texttt{cylinder\_penalty} & Rigid cylinder (penalty)\\
\hline
6 & \texttt{augmented\_Lagrangian\_KBC\_meshfree} & Kinematic constraints on meshfree nodes (augmented Lagrangian)\\
\hline
7 & \texttt{cylinder\_augmented\_Lagrangian} & Rigid cylinder (augmented Lagrangian)\\
\hline
8 & \texttt{field\_augmented\_Lagrangian\_KBC\_meshfree} & Field-driven kinematic constraints on meshfree nodes (augmented Lagrangian)\\
\hline
9 & \texttt{pressure\_bc} & Pressure boundary condition\\
\hline
10 & \texttt{angled\_bc} & Kinematic condition normal to an inclined surface (penalty)\\
\hline
11 & \texttt{penalty\_displacement\_meshfree} & Prescribed physical displacement on meshfree fields (penalty), Section~\ref{set.ICandBC.meshfree.EBC}\\
\hline
\end{tabular}}
\end{center}
\end{table}

\todo{This boundary condition type actually behaves more like 
``elements'' because they do produce a contribution to the stiffness
matrix, i.e., they act as ``follower forces'' and are not simply 
defined as functions of time.}\\

\subsubsection{Rigid barrier (penalty method)}
\subsubsection{Rigid sphere (penalty method)}
\subsubsection{Rigid sphere (augmented Lagrangian method)}
\subsubsection{Rigid sphere for meshfree domains (penalty method)}
\subsubsection{Rigid barrier (augmented Lagrangian method)}
\subsubsection{Rigid cylinder (penalty and augmented Lagrangian methods)}
\subsubsection{Kinematic constraints on meshfree nodes (augmented Lagrangian method)}
\subsubsection{Pressure boundary condition}
\subsubsection{Inclined kinematic condition (penalty method)}
\subsubsection{Prescribed physical displacement on meshfree fields (penalty method)}
See Section~\ref{set.ICandBC.meshfree.EBC}.

\subsection{Essential boundary conditions for meshfree fields}
\label{set.ICandBC.meshfree.EBC}
Meshfree shape functions do not satisfy $\Phi_{J}\farg{\mathbf{x}_{A}}=\delta_{JA}$
(Section~\ref{section.RKPM}), so the physical value at a node,
\begin{equation}
\label{eq.meshfree.EBC.value}
u^{h}\farg{\mathbf{x}_{A}} = \sum_{J} \Phi_{J}\farg{\mathbf{x}_{A}}\, d_{J},
\end{equation}
differs from its nodal coefficient $d_{A}$, and an ordinary kinematic boundary condition, which
prescribes $d_{A}$, does not impose the intended displacement. Two treatments are available. Both
obtain the values $\Phi_{J}\farg{\mathbf{x}_{A}}$ from the element group through a common
interface, which is currently implemented by the meshfree Kirchhoff--Love shell
(Section~\ref{section.meshfree.shell}).

\subsubsection{Direct nodal collocation}
The kinematic boundary condition controller \texttt{collocation\_KBC}, or equivalently the flag
\texttt{collocation="1"} on a \texttt{kinematic\_BC} entry, solves at every step for the
coefficients of the constrained nodes that make~\eqref{eq.meshfree.EBC.value} equal to the
prescribed value,
\begin{equation}
\mathbf{\Phi}_{c}\,\mathbf{d}_{c} = \bar{\mathbf{u}} - \mathbf{\Phi}_{f}\,\mathbf{d}_{f},
\end{equation}
where the subscripts $c$ and $f$ denote constrained and free coefficients, and imposes them as
ordinary prescribed values. For example,
\begin{inputfile}
<kinematic_BC dof="1" node_ID="3" schedule="1" type="u"
              value="-1.0" collocation="1"/>
\end{inputfile}
Collocation is exact for the prescribed value but is not a variational enforcement: the free
coefficients in the support of a constrained node receive no constraint force, and the reaction
reported for the constrained coefficient is not the physical load.

\subsubsection{Penalty enforcement of physical displacements}
The force controller \texttt{penalty\_displacement\_meshfree} enforces the constraint
$g_{A} = u^{h}\farg{\mathbf{x}_{A}} - \bar{u}_{A}\farg{t} = 0$ for one degree of freedom at every
node $A$ of the listed node sets by a penalty force distributed over the support,
\begin{equation}
f_{J} = \Phi_{J}\farg{\mathbf{x}_{A}}\,\lambda_{A}, \qquad
\lambda_{A} = -k\,g_{A} - c_{A}\,\dot{g}_{A}, \qquad
c_{A} = 2\,\zeta\,\sqrt{k\,m^{\textrm{eff}}_{A}}, \qquad
\frac{1}{m^{\textrm{eff}}_{A}} = \sum_{J}\frac{\Phi_{J}\farg{\mathbf{x}_{A}}^{2}}{m_{J}},
\end{equation}
where $k$ is the penalty stiffness, $\zeta$ the damping ratio of an optional dashpot that
suppresses oscillation of the penalty spring in explicit dynamics, and $m_{J}$ the lumped nodal
masses. The multiplier $\lambda_{A}$ is the physical reaction at node $A$, work conjugate to
$u^{h}\farg{\mathbf{x}_{A}}$. The controller contributes the stiffness
$k\,\Phi_{J}\Phi_{K}$, so it can also be used in implicit analyses. The parameters are
\texttt{dof}, \texttt{schedule}, \texttt{value} (the prescribed physical value, scaled by the
schedule), \texttt{penalty}, \texttt{damping\_ratio} (default 0) and \texttt{element\_group} (default 0,
which selects the first group that provides the shape function values). For example,
\begin{inputfile}
<penalty_displacement_meshfree dof="1" schedule="1"
    value="-300.0" penalty="1.0e5" damping_ratio="1.0">
  <node_ID_list><String value="3"/></node_ID_list>
</penalty_displacement_meshfree>
\end{inputfile}
At each output step the controller writes one line per constrained node to the output file with
the physical displacement, the reaction $\lambda_{A}$ and their means since the previous output.
The penalty stiffness should be large relative to the structural stiffness at the node but, in
explicit dynamics, small enough that $2\sqrt{k/m^{\textrm{eff}}_{A}}$ does not limit the stable
time step.


\subsection{Body forces}
\label{set.ICandBC.body}
Certain element formulations include body forces.
Applied body forces are specified with the vector analog 
of~\eqref{eq.func.time}.
Their values have the form
\begin{equation}
\label{eq.vec.func.time}
    \mathbf{b}\farg{t} = \mathbf{b}\,f^{\rbrkt{i}}\farg{t},
\end{equation}
where $\mathbf{b}$ is a vector coefficient of length $n_{dof}$, the 
number of degrees of freedom per node, and $f^{\rbrkt{i}}\farg{t}$ 
is a scalar schedule function. Schedule functions are explained in 
Section~\ref{sect.sched.param}.
Body forces are specified as
\[  
\sbrkt{\textit{schedule}} \quad 
\sbrkt{b_{1}} \quad \ldots \quad
\sbrkt{b_{n_{dof}}}
\]

\subsection{Natural boundary conditions}
\label{set.ICandBC.traction}
Recasting balance laws in weak form leads to natural boundary 
conditions involving flux terms across domain surfaces. These fluxes
are specified as part of the input of element parameters. Fluxes are 
associated with the faces of elements. As shown in Table~\ref{tab.BC.type}, 
the description of element faces differs depending on the geometry 
database format. Each natural boundary condition requires specification
of the flux vector at every node on the selected element face, or set 
of faces. Canonical number of element faces and node numbering over 
each face are described in Section~\ref{sect.integration.domains}.
Natural boundary conditions are specified as
\begin{align}
& \sbrkt{\textit{side ID}} \quad 
\sbrkt{\textit{schedule}} \quad 
\sbrkt{\textit{coordinate system}} \notag \\
& \begin{matrix}
\sbrkt{t_{11}} & \ldots & \sbrkt{t_{1n_{dof}}} \\
\vdots \\
\sbrkt{t_{n_{fn}1}} & \ldots & \sbrkt{t_{n_{fn}n_{dof}}}
\end{matrix} \notag
\end{align}
Each flux vector is length $n_{dof}$, and a vector must be specified 
for all $n_{fn}$ nodes on the element face.
The \textit{coordinate system} is identified by the values listed in 
Table~\ref{tab.BC.coord.sys}.
\begin{table}[h]
\begin{center}
\caption{\label{tab.BC.coord.sys} Coordinate systems for natural 
boundary conditions.}
\begin{tabular}[c]{|c|c|}
\hline
 \parbox[b]{0.75in}{\centering \textbf{code}}
&\parbox[b]{1.5in}{\centering \textbf{description}} \\
\hline
 \parbox[b]{0.75in}{\centering 0}
&\parbox[b]{1.5in}{\centering Cartesian} \\
\hline
 \parbox[b]{0.75in}{\centering 1}
&\parbox[b]{1.5in}{\centering facet local} \\
\hline
\end{tabular}
\end{center}
\end{table}
The facet local coordinate frame is constructed such that the last 
component of the flux vector $t_{n_{dof}}$ corresponds to flux in the 
outward normal direction. The tangential components of the flux vector a 
constructed from the mapping of the face to its corresponding parent 
integration domain.

An example specification of two-dimensional, traction boundary conditions appears below:
\begin{inputfile}
1	# number of traction boundary conditions
####
1  # side ID
1  # schedule number
0  # coordinate system 
# traction vectors
0.0  1.0   
0.0  1.0
\end{inputfile}
The example shows a single specification. As shown in Table~\ref{tab.BC.type}, the format of the \textit{side ID} specification depends on the geometry database format. In the example, a single integer is used to define the element edges where the boundary conditions are applied. This format would be appropriate for geometry database types that support definition of side sets. Two traction vectors, each with two components, are given indicating that the element edges specified by the \textit{side ID} have two nodes per edge. The traction vectors are specified with respect to the global Cartesian coordinate frame, meaning the values for each vector correspond to $\cbrkt{t_x, t_y}$. A traction vector must be provided for each node of the element face. The number of nodes on each face for a given element geometry and the canonical order of the nodes on each face is described in Section~\ref{sect.integration.domains}. The traction is interpolated over the face from the nodal values using the standard finite element shape functions.

